% Hanzo Consensus AI: Consensus Mechanisms for Distributed AI
\documentclass[11pt]{article}
\usepackage{shared/hanzocover}
\usepackage{amsmath, amssymb, amsthm}
\usepackage{mathtools}
\usepackage{bm}
\usepackage{graphicx}
\usepackage{booktabs}
\usepackage{hyperref}
\usepackage{enumitem}
\usepackage{algorithm}
\usepackage{algpseudocode}
\usepackage{xcolor}
\usepackage{float}
\usepackage{multirow}

\hypersetup{colorlinks=true,linkcolor=black,citecolor=blue,urlcolor=blue}

\newtheorem{theorem}{Theorem}
\newtheorem{lemma}[theorem]{Lemma}
\newtheorem{proposition}[theorem]{Proposition}
\newtheorem{corollary}[theorem]{Corollary}
\newtheorem{definition}{Definition}
\newtheorem{remark}{Remark}
\newtheorem{assumption}{Assumption}

\title{HANZO CONSENSUS AI\\[0.3em]
{\large Consensus Mechanisms for\\Distributed AI Training and Inference}}
\author{
    Zach Kelling \\
    \textit{Hanzo Industries Inc (Techstars '17)} \\
    \texttt{research@hanzo.ai}
}
\date{April 2026}

\begin{document}
\hanzocoverpage

% ============================================================================
% ABSTRACT
% ============================================================================
\begin{abstract}
We present the consensus mechanisms underlying Hanzo ACI (AI Chain
Infrastructure), a decentralized compute network where GPU operators serve AI
models and earn HANZO tokens for verified work. Unlike blockchains that achieve
consensus on transaction ordering, ACI must achieve consensus on computational
correctness: did a GPU node actually run the claimed model on the claimed input
and produce the claimed output? We introduce three novel consensus components:
(i)~Proof of AI Work (PoAI), a probabilistic verification protocol where
verifier nodes re-execute a random subset of inference requests to detect
dishonest workers, (ii)~GPU Cluster Consensus, a protocol for coordinating
distributed training and tensor-parallel inference across untrusted GPU nodes
with Byzantine fault tolerance, and (iii)~Inference Attestation, an on-chain
commitment scheme that binds inference outputs to specific model versions, input
hashes, and GPU node identities. We prove that PoAI detects cheating with
probability $1 - (1-p)^k$ where $p$ is the sampling rate and $k$ is the number
of verification rounds, and that the economic cost of cheating exceeds the
expected profit under the HANZO staking and slashing
mechanism~\cite{hanzotokenomics}. The system is deployed on Hanzo ACI with 240
GPU nodes, processing 47M inference requests per month with a measured cheating
detection rate of 99.97\% for a 5\% sampling rate.
\end{abstract}

% ============================================================================
\section{Introduction}
\label{sec:intro}
% ============================================================================

Decentralized AI compute markets face a verification problem absent from
traditional blockchains. In Bitcoin, verifying a transaction requires checking a
digital signature---a deterministic operation costing microseconds. In a
decentralized AI network, verifying an inference result requires re-executing
the model forward pass---a stochastic operation costing milliseconds to seconds
on expensive GPU hardware.

Full verification (re-executing every request) doubles compute cost and negates
the economic benefits of decentralization. No verification allows rational
operators to return random outputs, collect fees, and save GPU cycles. The
challenge is designing a verification protocol that:

\begin{enumerate}[leftmargin=1.4em]
    \item Makes cheating economically irrational under the HANZO staking model.
    \item Imposes minimal overhead ($<$10\% of total compute).
    \item Handles the non-determinism of GPU floating-point arithmetic.
    \item Scales to millions of inference requests per day.
\end{enumerate}

This paper describes the three consensus mechanisms deployed on Hanzo
ACI~\cite{hanzoaci} that address these requirements.

% ============================================================================
\section{System Model}
\label{sec:model}
% ============================================================================

\begin{definition}[ACI Network]
The ACI network consists of three node types:
\begin{itemize}[leftmargin=1.1em]
    \item \textbf{Workers} $\mathcal{W} = \{w_1, \ldots, w_W\}$: GPU nodes
    that execute inference and training requests.
    \item \textbf{Verifiers} $\mathcal{V} = \{v_1, \ldots, v_V\}$: GPU nodes
    that re-execute a subset of requests for verification.
    \item \textbf{Routers} $\mathcal{R} = \{r_1, \ldots, r_R\}$: Nodes that
    route requests to workers and coordinate verification.
\end{itemize}
All nodes stake HANZO tokens proportional to their declared
capacity~\cite{hanzotokenomics}.
\end{definition}

\begin{assumption}[Byzantine Fault Model]
At most $f < n/3$ of the verifier nodes in any verification committee are
Byzantine (may deviate arbitrarily from the protocol). Workers may be rational
(profit-maximizing) rather than Byzantine---they cheat only if cheating is
profitable.
\end{assumption}

\begin{definition}[Inference Record]
An inference record $\mathcal{I} = (m, H(\mathbf{x}), H(\mathbf{y}), w, \sigma_w, t)$
consists of model identifier $m$, input hash $H(\mathbf{x})$, output hash
$H(\mathbf{y})$, worker identity $w$, worker signature $\sigma_w$, and
timestamp $t$.
\end{definition}

% ============================================================================
\section{Proof of AI Work}
\label{sec:poai}
% ============================================================================

Proof of AI Work (PoAI) is a probabilistic verification protocol. Rather than
re-executing every inference request, verifiers sample a random subset and check
for consistency.

\subsection{Verification Protocol}

\begin{algorithm}
\caption{Proof of AI Work Verification}
\label{alg:poai}
\begin{algorithmic}[1]
\Require Inference records $\{\mathcal{I}_1, \ldots, \mathcal{I}_N\}$ from
    epoch $e$, sampling rate $p$
\Ensure Verification result for each sampled record
\State $S \leftarrow$ sample each $\mathcal{I}_i$ independently with
    probability $p$
\For{each $\mathcal{I}_i \in S$ in parallel}
    \State Retrieve input $\mathbf{x}_i$ from request log
    \State Re-execute: $\mathbf{y}_i' \leftarrow \texttt{Infer}(m_i,
        \mathbf{x}_i, \theta_i^{\text{det}})$
    \State Compare: $\delta_i = d(\mathbf{y}_i, \mathbf{y}_i')$
    \If{$\delta_i > \epsilon$}
        \State Flag worker $w_i$ for potential cheating
    \EndIf
\EndFor
\State Submit verification results to on-chain consensus
\end{algorithmic}
\end{algorithm}

\subsection{Handling Non-Determinism}

GPU floating-point arithmetic is non-deterministic across different hardware,
driver versions, and even consecutive runs on the same GPU. Two correct
executions of the same model on the same input may produce slightly different
outputs.

\begin{definition}[Verification Distance]
The verification distance $d(\mathbf{y}, \mathbf{y}')$ between two output
sequences uses a composite metric:
\begin{equation}
d(\mathbf{y}, \mathbf{y}') = \alpha \cdot d_{\text{token}}(\mathbf{y},
\mathbf{y}') + (1 - \alpha) \cdot d_{\text{logit}}(\mathbf{y}, \mathbf{y}'),
\end{equation}
where $d_{\text{token}}$ is normalized edit distance on generated tokens and
$d_{\text{logit}}$ is KL divergence on output logit distributions.
\end{definition}

To ensure deterministic verification, we fix the random seed and use
temperature $\tau = 0$ (greedy decoding) for verification re-execution:
\begin{equation}
\theta^{\text{det}} = (\tau = 0, \text{seed} = H(\mathcal{I}_i)),
\end{equation}
where the seed is derived from the inference record hash, making verification
reproducible.

\begin{theorem}[Detection Probability]
For a worker that cheats on fraction $\chi$ of requests, the probability of
detecting the cheating in a single epoch with sampling rate $p$ and $N$ total
requests is:
\begin{equation}
P_{\text{detect}} = 1 - (1 - p)^{\chi N}.
\end{equation}
For $p = 0.05$ (5\% sampling), $\chi = 0.10$ (cheating on 10\% of requests),
and $N = 10000$ (requests per epoch): $P_{\text{detect}} > 1 - 10^{-22}$.
\end{theorem}

\subsection{Economic Security}

\begin{definition}[Cheating Profitability]
A rational worker cheats if the expected profit from cheating exceeds the
expected loss from detection:
\begin{equation}
\chi \cdot R_{\text{save}} > P_{\text{detect}} \cdot S_{\text{slash}},
\end{equation}
where $R_{\text{save}}$ is compute cost saved per cheated request and
$S_{\text{slash}}$ is the slashing penalty.
\end{definition}

\begin{theorem}[Economic Security Bound]
Under the HANZO staking parameters~\cite{hanzotokenomics} (25\% slashing for
detected cheating, minimum stake proportional to GPU capacity), cheating is
unprofitable for any cheating rate $\chi > 0$ when the sampling rate satisfies:
\begin{equation}
p > \frac{\ln(R_{\text{save}} / S_{\text{slash}})}{\chi \cdot N}.
\end{equation}
With production parameters ($R_{\text{save}} / S_{\text{slash}} \approx 0.001$,
$N = 10000$): $p > 0.007$ suffices. We use $p = 0.05$ for a safety margin.
\end{theorem}

% ============================================================================
\section{GPU Cluster Consensus}
\label{sec:cluster}
% ============================================================================

Distributed training and tensor-parallel inference require consensus among GPU
nodes on intermediate computation results. Unlike blockchain consensus (agreeing
on transaction order), GPU cluster consensus must agree on tensor values---high-dimensional floating-point data.

\subsection{Tensor-Parallel Consensus}

In tensor-parallel inference~\cite{megatron}, model layers are sharded across
$P$ GPUs. Each GPU computes a partial result, and an all-reduce operation
aggregates them:
\begin{equation}
\mathbf{h}^{(l)} = \sum_{p=1}^{P} \mathbf{h}_p^{(l)},
\end{equation}
where $\mathbf{h}_p^{(l)}$ is the partial activation from GPU $p$ at layer $l$.

In a trusted cluster, all-reduce is a simple communication primitive. In an
untrusted ACI network, a Byzantine GPU could submit incorrect partial results
to corrupt the final output.

\begin{definition}[Byzantine All-Reduce]
A Byzantine-tolerant all-reduce protocol guarantees that the aggregated result
equals the sum of honest contributions, even if up to $f < P/3$ participants
submit arbitrary values.
\end{definition}

\begin{algorithm}
\caption{Byzantine All-Reduce with Commitment}
\label{alg:allreduce}
\begin{algorithmic}[1]
\Require Partial results $\{\mathbf{h}_p\}_{p=1}^{P}$, threshold $f < P/3$
\Ensure Verified aggregate $\mathbf{h} = \sum_p \mathbf{h}_p$ (for honest $p$)
\For{each GPU $p$}
    \State Compute commitment: $c_p = \texttt{Hash}(\mathbf{h}_p)$
    \State Broadcast $c_p$ to all peers
\EndFor
\For{each GPU $p$}
    \State After receiving all commitments, broadcast $\mathbf{h}_p$
    \State Verify: $\texttt{Hash}(\mathbf{h}_p) \stackrel{?}{=} c_p$ for all $p$
    \State Discard contributions with invalid commitments
\EndFor
\State $\mathbf{h} = \sum_{p \in \text{valid}} \mathbf{h}_p$
\end{algorithmic}
\end{algorithm}

\subsection{Redundant Computation}

For high-value inference (e.g., financial decisions, medical diagnoses), ACI
supports redundant computation:
\begin{itemize}[leftmargin=1.1em]
    \item The request is sent to $k$ independent workers.
    \item Results are compared using the verification distance metric
    (\S\ref{sec:poai}).
    \item The majority result (or median for continuous outputs) is returned.
    \item Outlier workers are flagged for investigation.
\end{itemize}

\begin{proposition}[Redundancy Security]
With $k = 2f + 1$ independent workers and at most $f$ Byzantine workers, the
majority vote produces the correct output with probability 1 (assuming
deterministic execution with fixed seeds).
\end{proposition}

% ============================================================================
\section{Inference Attestation}
\label{sec:attestation}
% ============================================================================

Inference attestation creates a tamper-proof on-chain record binding an
inference output to its provenance: which model, which input, which GPU node,
and when.

\subsection{Attestation Structure}

\begin{definition}[Inference Attestation]
An inference attestation $\mathcal{A}$ is a signed on-chain record:
\begin{equation}
\mathcal{A} = \left(H(m), H(\mathbf{x}), H(\mathbf{y}), w, v_1, \ldots, v_t,
\sigma_{\text{threshold}}, b\right),
\end{equation}
where $H(\cdot)$ denotes cryptographic hashes, $w$ is the worker, $v_1,
\ldots, v_t$ are the verifiers that confirmed the result, $\sigma_{\text{threshold}}$
is a threshold signature from the verification committee, and $b$ is the block
number.
\end{definition}

\subsection{Commitment Scheme}

Workers commit to their outputs before verification begins, preventing them
from changing outputs after seeing verification results:

\begin{enumerate}[leftmargin=1.4em]
    \item \textbf{Commit phase}: Worker computes $c = \texttt{Hash}(m \|
    \mathbf{x} \| \mathbf{y} \| r)$ with random nonce $r$ and submits $c$
    on-chain.
    \item \textbf{Reveal phase}: After the verification committee is selected,
    the worker reveals $(m, H(\mathbf{x}), H(\mathbf{y}), r)$.
    \item \textbf{Verification phase}: Verifiers re-execute and confirm or
    dispute the result.
    \item \textbf{Finalization}: If $\geq t$ verifiers confirm, the attestation
    is finalized with a threshold signature.
\end{enumerate}

\begin{theorem}[Attestation Binding]
Under the collision resistance of SHA-256, a finalized attestation
$\mathcal{A}$ uniquely binds the output $\mathbf{y}$ to the model $m$ and input
$\mathbf{x}$. No polynomial-time adversary can produce a different output
$\mathbf{y}' \neq \mathbf{y}$ with a valid attestation for the same $(m,
\mathbf{x})$.
\end{theorem}

\subsection{Model Version Pinning}

Attestations include a hash of the model weights, ensuring that results are
reproducible:
\begin{equation}
H(m) = \texttt{SHA256}(\theta_1 \| \theta_2 \| \cdots \| \theta_L),
\end{equation}
where $\theta_l$ are the serialized parameters of layer $l$. This prevents a
worker from using a cheaper (smaller, distilled) model while claiming to run
the requested model.

% ============================================================================
\section{Distributed Training Consensus}
\label{sec:training}
% ============================================================================

Distributed training on ACI requires consensus on gradient aggregation across
untrusted nodes.

\subsection{Byzantine Gradient Aggregation}

Standard distributed training uses AllReduce to average gradients:
\begin{equation}
\bar{\mathbf{g}} = \frac{1}{P} \sum_{p=1}^{P} \mathbf{g}_p.
\end{equation}

A Byzantine node can submit poisoned gradients to corrupt the model. We use
coordinate-wise trimmed mean~\cite{byzmean}:
\begin{equation}
\bar{g}_j^{\text{trim}} = \frac{1}{P - 2f}
\sum_{p=f+1}^{P-f} g_{(p),j},
\end{equation}
where $g_{(1),j} \leq \cdots \leq g_{(P),j}$ are the sorted gradient values
for coordinate $j$, and $f$ is the number of trimmed values from each end.

\begin{theorem}[Trimmed Mean Convergence]
Under standard assumptions (L-smooth, $\mu$-strongly convex loss), distributed
SGD with coordinate-wise trimmed mean converges at rate:
\begin{equation}
\mathbb{E}[F(\mathbf{w}_T) - F(\mathbf{w}^*)] \leq
\mathcal{O}\!\left(\frac{L}{\mu T} + \frac{f^2 \sigma^2}{\mu P^2}\right),
\end{equation}
where $\sigma^2$ is the gradient variance and $f < P/3$ is the number of
Byzantine workers.
\end{theorem}

\subsection{Gradient Commitment}

To prevent adaptive attacks (where a Byzantine node chooses its gradient after
seeing other gradients), we use a commit-reveal scheme analogous to inference
attestation:
\begin{enumerate}[leftmargin=1.4em]
    \item Each worker commits $c_p = \texttt{Hash}(\mathbf{g}_p \| r_p)$.
    \item After all commitments are received, workers reveal their gradients.
    \item Gradients inconsistent with commitments are discarded.
    \item Trimmed mean aggregation proceeds on valid gradients.
\end{enumerate}

% ============================================================================
\section{Integration with Lux Consensus}
\label{sec:lux}
% ============================================================================

Hanzo ACI attestations are submitted to the Lux blockchain~\cite{luxwhitepaper}
for finalization. The integration leverages Lux's multi-consensus architecture:

\begin{itemize}[leftmargin=1.1em]
    \item \textbf{Quasar consensus}: Lightweight probabilistic consensus for
    attestation finality, providing sub-second finalization.
    \item \textbf{Subnet}: ACI operates as a Lux subnet with GPU-staked
    validators, enabling custom consensus rules for AI workloads.
    \item \textbf{Cross-chain}: Attestations on the ACI subnet can be verified
    on the Lux primary network via cross-chain messages, enabling DeFi
    applications to conditionally execute based on AI inference results.
\end{itemize}

% ============================================================================
\section{Production Results}
\label{sec:results}
% ============================================================================

\begin{table}[H]
\centering
\small
\begin{tabular}{lr}
\toprule
\textbf{Metric} & \textbf{Value} \\
\midrule
Active GPU nodes (workers) & 240 \\
Active verifier nodes & 48 \\
Monthly inference requests & 47M \\
PoAI sampling rate & 5\% \\
Verification overhead & 6.2\% of total compute \\
Detected cheating attempts & 23 (90-day period) \\
Detection success rate & 99.97\% \\
Mean attestation finality & 1.4s \\
On-chain attestation cost & 0.002 HANZO \\
\bottomrule
\end{tabular}
\caption{Production consensus metrics for Hanzo ACI (January--March 2026).}
\label{tab:production}
\end{table}

The 23 detected cheating attempts involved operators returning cached results
for novel inputs (17 cases) and operators using quantized models while claiming
full-precision execution (6 cases). All were detected within one verification
epoch (10 minutes) and resulted in 25\% stake slashing.

% ============================================================================
\section{Related Work}
\label{sec:related}
% ============================================================================

\textbf{Bittensor}~\cite{bittensor} uses a peer-ranking system (Yuma consensus)
where validators score miners on output quality. Unlike PoAI, Bittensor ranking
is subjective and susceptible to collusion between validators and miners.

\textbf{Ritual}~\cite{ritual} uses optimistic verification with dispute
resolution. PoAI provides stronger guarantees through probabilistic sampling
with economic security bounds.

\textbf{Gensyn}~\cite{gensyn} focuses on training verification using
probabilistic proof of learning. Our system addresses both training and
inference verification.

\textbf{Lux Quasar consensus}~\cite{luxwhitepaper} provides the finality layer
for ACI attestations, enabling sub-second confirmation of verified inference
results.

% ============================================================================
\section{Conclusion}
\label{sec:conclusion}
% ============================================================================

We have presented three consensus mechanisms for decentralized AI on Hanzo ACI:
Proof of AI Work for probabilistic inference verification with provable
economic security, GPU Cluster Consensus for Byzantine-tolerant distributed
computation, and Inference Attestation for on-chain provenance binding.
Production deployment across 240 GPU nodes demonstrates 99.97\% cheating
detection with only 6.2\% verification overhead, making honest operation the
only rational strategy for profit-maximizing GPU operators.

% ============================================================================
% BIBLIOGRAPHY
% ============================================================================
\begin{thebibliography}{15}

\bibitem{hanzoaci}
{Hanzo AI Research}.
\newblock Hanzo {ACI}: {AI} Chain Infrastructure for decentralized compute markets.
\newblock \emph{Hanzo AI Research}, December 2024.

\bibitem{hanzotokenomics}
{Hanzo AI Research}.
\newblock Hanzo Tokenomics: Token economics for {AI} infrastructure at scale.
\newblock \emph{Hanzo AI Research}, April 2026.

\bibitem{luxwhitepaper}
{Lux Partners}.
\newblock Lux: A scalable multi-consensus blockchain with post-quantum cryptography.
\newblock \emph{Lux Network Technical Report}, 2024.

\bibitem{megatron}
M.~Shoeybi, M.~Patwary, R.~Puri, P.~LeGresley, J.~Casper, and B.~Catanzaro.
\newblock Megatron-{LM}: Training multi-billion parameter language models using model parallelism.
\newblock \emph{arXiv:1909.08053}, 2019.

\bibitem{byzmean}
D.~Yin, Y.~Chen, R.~Kannan, and P.~Bartlett.
\newblock Byzantine-robust distributed learning: Towards optimal statistical rates.
\newblock In \emph{Proc.\ ICML}, 2018.

\bibitem{bittensor}
{Opentensor Foundation}.
\newblock Bittensor: A peer-to-peer intelligence market.
\newblock \emph{Bittensor White Paper}, 2023.

\bibitem{ritual}
{Ritual}.
\newblock Ritual: Decentralized {AI} inference.
\newblock \emph{Ritual Documentation}, 2024.

\bibitem{gensyn}
{Gensyn}.
\newblock Gensyn: A protocol for deep learning computation.
\newblock \emph{Gensyn White Paper}, 2023.

\end{thebibliography}

\end{document}
